% Research synthesis of two archived held-out graph studies, 10 Nov 2026.
% All values in tables below originate from research-branch
% M5/M6 paired_contrasts.csv, snapshot_summary.csv and quality_summary.csv.
\section{Comparative interpretation: numerical fidelity, optimizer quality, and cost}
\label{sec:heldout-synthesis}

\subsection{A graph-conditioned reversal already occurs in the local field}

The two held-out instances both have 100 spins, but are
not interchangeable workloads: M5 \texttt{ising2.5-100\_7777}
has 4,869 nonzero couplings among 4,950 original source entries,
whereas M6 \texttt{t2g10\_7777} has 200 nonzero couplings.
In each study the original graph weights determine a
source-only global scale \emph{separately} for FP6 and INT6.
The graph-dependent fitted scale is part of the declared
representation, rather than a parameter tuned using Max-Cut
solution quality.

\begin{table}[htbp]\centering\small
\begin{tabular}{@{}c rrrr@{}}\toprule
 & \multicolumn{2}{c}{Dense M5} & \multicolumn{2}{c}{Sparse M6}\\
$d$ & FP6 (\%) & INT6 (\%) & FP6 (\%) & INT6 (\%)\\\midrule
0.0 & 3.501 & 4.835 & 3.414 & 1.875\\
0.8 & 3.516 & 4.786 & 3.436 & 1.906\\
1.6 & 3.538 & 4.790 & 3.406 & 1.900\\\bottomrule
\end{tabular}
\caption{Mean state-conditioned relative RMS local-field error
(relative fraction times 100), extracted from each archived
\texttt{snapshot\_summary.csv}. At all three assumed offset
levels, FP6 has lower field error on dense M5 but INT6
has lower field error on sparse M6. Values describe
common FP64-returned reference states, not dynamic
solver trajectories.}
\label{tab:heldout-relativefield}
\end{table}

The reversal in static switching-probability error therefore
cannot be attributed \emph{solely} to the sigmoid output
nonlinearity: the ranking of the sampled local-field
errors has also changed. Plausible but \emph{untested}
explanations include the graph weight-magnitude distribution,
the source-only globally fitted scales, saturation and
zeroing of coefficients, local-field cancellation, and
different state distributions. Only two individual
instances were studied; topology, weight distribution
and visited states were not isolated as experimental
factors. We cannot claim a causal sparsity effect.

\subsection{The two frozen quality screens answer a different question}

Define the extra-gap screen relative to the same-response
FP64 baseline. The following are the \emph{upper limits}
of the paired seed-bootstrap 95\% confidence intervals;
these are not the estimated mean extra gaps.
\begin{table}[htbp]\centering\small
\begin{tabular}{@{}c rr rr@{}}\toprule
 & \multicolumn{2}{c}{Dense M5 (pp)}
 & \multicolumn{2}{c}{Sparse M6 (pp)}\\
$d$ & FP6 & INT6 & FP6 & INT6\\\midrule
0.0 & $+0.478$ & $+1.245$ & $+1.522$ & $+1.121$\\
0.8 & $+0.241$ & $+1.287$ & $+0.278$ & $+0.677$\\
1.6 & $+0.716$ & $+0.952$ & $+0.717$ & $+1.141$\\\bottomrule
\end{tabular}
\caption{Paired additional optimum-gap 95\% CI upper
bounds against same-offset FP64 from the archived
\texttt{quality\_summary.csv}. Eligibility additionally
requires the \emph{absolute} mean gap upper CI not
to exceed 5\%. Passing the +1-pp test alone is
insufficient: dense M5 FP6 at $d=1.6$ fails the
absolute screen with upper bound 5.286\%.}
\label{tab:heldout-extra-ci}
\end{table}

On sparse M6, both FP6 and INT6 pass \emph{both}
predeclared screens at $d=0.8$, and neither encoding
has a statistically certified \emph{direct} advantage
over the other in final original-weight Max-Cut gap:
the paired FP6-minus-INT6 CI is
$[-1.650,+0.034]$ pp. At $d=1.6$, sparse FP6 passes
the two engineering screens and INT6 does not, even
though INT6 has lower state-conditioned probability
distortion. The direct cut-gap CI still crosses zero,
$[-1.378,+0.386]$ pp. This difference in
threshold certification is not, by itself, proof of
an inferential format-ranking difference.

For M6 the frozen within-graph offset-by-format
difference-in-differences (FP6-minus-INT6 gap
contrast against $d=0$) is $-1.167$ pp with
95\% CI $[-2.791,+0.468]$ at $d=0.8$ and
$-0.845$ pp with CI $[-2.215,+0.556]$ at $d=1.6$.
Both CIs span zero. The corresponding M5 interactions
also span zero. No format-by-offset interaction has
been demonstrated; a formal across-graph interaction
comparison has not been performed. An interval
containing zero is inconclusive, not an equivalence test.

\subsection{Implication for hardware-aware representation selection}

At equal six-bit precision, FP6 and INT6 have identical
nominal payload bits per stored undirected edge, plus
one binary32 source-only global scale in the existing
software representation. The existing 99-term generic
mapped digital row-plus-interface comparison
measures $23481.2704\,\mu\mathrm{m}^2$ for FP6
versus $18061.0720\,\mu\mathrm{m}^2$ for INT6,
a ratio of approximately $1.30$. This is
architecture- and library-specific standard-cell
area, \emph{not} total processor area, memory
area, energy per solution, or a measured sparse
graph implementation. The simulation uses cached
incremental local fields, not the same serial
recomputation as the mapped example.

The current evidence supports \emph{conditional}
format choice rather than universal FP preference.
Static field/probability fidelity should be measured,
but the frozen finite-budget quality target must be
evaluated separately. A future FP8-versus-INT8
comparison on the already-viewed dense graph is
a distinct proposed study and has \emph{not}
been included in the M5/M6 results.

\noindent\textbf{Reproducibility and scope.}
Both original 216-run experimental archives,
their protocol/graph SHA256 manifests, complete
paired uncertainty summaries and figures are
retained on the \href{https://github.com/Makise-okabe/MRAM_P-bit/tree/research/m4-gain-fusion-20261107/results}{research branch}.
Each graph has a \emph{single} persistent synthetic
device-offset population; seed-bootstrap confidence
intervals cannot be construed as manufacturing
statistics or graph-population confidence intervals.
The present STT-MRAM p-bit switching-response
variations are assumed E2 stresses, not measured
post-calibration physical residuals.
